我们提出 \wcabench{}——首个基于世界魔方协会（\wca）全部公开比赛数据的机器学习标准化基准。速拧是一种独特而富有启发性的竞技场景：它记录了延续数十年、覆盖全球、以个人为单位的\emph{数百万}次计时尝试，并受一套显式规则体系约束（轮次格式、“5 次去极值平均 \code{ao5}”、多盲复合编码、DNF/DNS 状态码），而通用时序与表格管线并不建模这些规则。基于包含 \textbf{6{,}909{,}454} 条成绩、\textbf{298{,}551} 名选手与 \textbf{18{,}708} 场比赛的原始导出，我们定义了时间划分（train 2003--2022、val 2023--2024、test 2025--2026）以及覆盖回归、排序、不平衡分类、极值外推与因果推断的五类任务：\task{1} 成绩预测、\task{2} 名次预测、\task{3} DNF 预测、\task{4} 人类极限估计与 \task{5} 技能迁移。本文有四项贡献：(i) 可复现的列式数据基础设施，含规则感知解码与数据卡；(ii) 为五个任务给出形式化、指标完备的定义并刻画领域挑战；(iii) 防泄漏的滚动窗口评估协议，含冻结统计量、四维分层报告与完整的统计检验工具箱；(iv) 覆盖六大方法家族（统计、树模型、深度、图、贝叶斯与因果）的 \textbf{21 个参考基线}及真实实测数字——例如基于树模型的成绩预测（单次运行）\maelog{} $=0.164$、官方 Psych Sheet 名次预测 Kendall $\kendall=0.767$、基于树模型的 DNF 预测 \aucpr{} $=0.373$。我们还如实报告了简单统计模型仍然具有竞争力、以及外推与混淆会使结论不可靠之处。
